pdfoutput=1
% Options for packages loaded elsewhere
\PassOptionsToPackage{unicode}{hyperref}
\PassOptionsToPackage{hyphens}{url}
\pdfoutput=1
\documentclass[
  12pt,
]{article}
\usepackage{xcolor}
\usepackage{amsmath,amssymb}
\setcounter{secnumdepth}{-\maxdimen} % remove section numbering
\usepackage{iftex}
\ifPDFTeX
  \usepackage[T1]{fontenc}
  \usepackage{textcomp} % provide euro and other symbols
\else % if luatex or xetex
  \usepackage{unicode-math} % this also loads fontspec
  \defaultfontfeatures{Scale=MatchLowercase}
  \defaultfontfeatures[\rmfamily]{Ligatures=TeX,Scale=1}
\fi
\usepackage{lmodern}
\ifPDFTeX\else
  % xetex/luatex font selection
\fi
% Use upquote if available, for straight quotes in verbatim environments
\IfFileExists{upquote.sty}{\usepackage{upquote}}{}
\IfFileExists{microtype.sty}{% use microtype if available
  \usepackage[]{microtype}
  \UseMicrotypeSet[protrusion]{basicmath} % disable protrusion for tt fonts
}{}
\makeatletter
\@ifundefined{KOMAClassName}{% if non-KOMA class
  \IfFileExists{parskip.sty}{%
    \usepackage{parskip}
  }{% else
    \setlength{\parindent}{0pt}
    \setlength{\parskip}{6pt plus 2pt minus 1pt}}
}{% if KOMA class
  \KOMAoptions{parskip=half}}
\makeatother
\usepackage{longtable,booktabs,array}
\usepackage{calc} % for calculating minipage widths
% Correct order of tables after \paragraph or \subparagraph
\usepackage{etoolbox}
\makeatletter
\patchcmd\longtable{\par}{\if@noskipsec\mbox{}\fi\par}{}{}
\makeatother
% Allow footnotes in longtable head/foot
\IfFileExists{footnotehyper.sty}{\usepackage{footnotehyper}}{\usepackage{footnote}}
\makesavenoteenv{longtable}
\setlength{\emergencystretch}{3em} % prevent overfull lines
\providecommand{\tightlist}{%
  \setlength{\itemsep}{0pt}\setlength{\parskip}{0pt}}
\usepackage{bookmark}
\IfFileExists{xurl.sty}{\usepackage{xurl}}{} % add URL line breaks if available
\urlstyle{same}
\hypersetup{
  hidelinks,
  pdfcreator={LaTeX via pandoc}}

\author{SCX}
\date{2026}

\begin{document}

\section{Situs Theory Final Verification Report}\label{situs-theory-final-verification-report}

\textbf{Verification Date}: 2026-06-29\\
\textbf{Verified Files}: - \texttt{paper/situs\_theory/main.tex} (1518 lines) - \texttt{theory/self\_evolution/ppe\_rigorous\_derivation.md} (1030 lines)

\begin{center}\rule{0.5\linewidth}{0.5pt}\end{center}

\subsection{Item-by-Item Check}\label{item-by-item-check}

\subsubsection{1. Theorem 1.2.3 Lipschitz Constant: Corrected to $2\sqrt{2}\cdot\pi/\sqrt{d}$?}\label{thm-1.2.3-lipschitz-check}

\begin{longtable}[]{@{}
  >{\raggedright\arraybackslash}p{(\linewidth - 6\tabcolsep) * \real{0.2500}}
  >{\raggedright\arraybackslash}p{(\linewidth - 6\tabcolsep) * \real{0.2500}}
  >{\raggedright\arraybackslash}p{(\linewidth - 6\tabcolsep) * \real{0.2500}}
  >{\raggedright\arraybackslash}p{(\linewidth - 6\tabcolsep) * \real{0.2500}}@{}}
\toprule\noalign{}
\begin{minipage}[b]{\linewidth}\raggedright
File
\end{minipage} & \begin{minipage}[b]{\linewidth}\raggedright
Location
\end{minipage} & \begin{minipage}[b]{\linewidth}\raggedright
Formula
\end{minipage} & \begin{minipage}[b]{\linewidth}\raggedright
Status
\end{minipage} \\
\midrule\noalign{}
\endhead
\bottomrule\noalign{}
\endlastfoot
main.tex & L405--409 & \texttt{L\_\{{\textbackslash}PE\}$\hat{\ }$\{{\textbackslash}text\{scalar\}\} = {\textbackslash}frac\{2{\textbackslash}sqrt\{2\}{\textbackslash},{\textbackslash}pi\}\{{\textbackslash}sqrt\{d\}\} {\textbackslash}cdot {\textbackslash}sqrt\{{\textbackslash}sum\_\{j=0\}$\hat{\ }$\{d/2-1\}{\textbackslash}frac\{1\}\{{\textbackslash}lambda\_j$\hat{\ }$2\}\}} & \checkmark \\
ppe\_rigorous\_derivation.md & L189 & Same form & \checkmark \\
\end{longtable}

\textbf{Conclusion}: Both files consistent, constant corrected. The normalization factor $\sqrt{2/d}$ has been incorporated into the proof derivation, with the final Lipschitz constant obtaining coefficient $2\sqrt{2}\cdot\pi/\sqrt{d}$ through the sum over $\sqrt{2/d} \times 2\pi/\lambda_j$. The asymptotic behavior of Corollary 1.2.3 $\sim (1/\xi)\cdot\sqrt{d/6}$ consequently drops from $O(d)$ to $O(\sqrt{d})$. \textbf{PASS.}

\begin{center}\rule{0.5\linewidth}{0.5pt}\end{center}

\subsubsection{2. Theorem 1.2.1: Normalization Factor Added? $L$-dependence removed from boxed formulas?}\label{thm-1.2.1-normalization-check}

\textbf{2a. Normalization factor $\sqrt{2/d}$}:

\begin{longtable}[]{@{}
  >{\raggedright\arraybackslash}p{(\linewidth - 6\tabcolsep) * \real{0.2500}}
  >{\raggedright\arraybackslash}p{(\linewidth - 6\tabcolsep) * \real{0.2500}}
  >{\raggedright\arraybackslash}p{(\linewidth - 6\tabcolsep) * \real{0.2500}}
  >{\raggedright\arraybackslash}p{(\linewidth - 6\tabcolsep) * \real{0.2500}}@{}}
\toprule\noalign{}
\begin{minipage}[b]{\linewidth}\raggedright
File
\end{minipage} & \begin{minipage}[b]{\linewidth}\raggedright
Location
\end{minipage} & \begin{minipage}[b]{\linewidth}\raggedright
Content
\end{minipage} & \begin{minipage}[b]{\linewidth}\raggedright
Status
\end{minipage} \\
\midrule\noalign{}
\endhead
\bottomrule\noalign{}
\endlastfoot
main.tex & L268--269 & $\mathsf{PE}_{\text{scalar}}(p, 2j) = \sqrt{\frac{2}{d}}\,\sin(\ldots)$ & \checkmark \\
ppe\_rigorous\_derivation.md & L101 & Same form & \checkmark \\
\end{longtable}

Both files include the normalization factor and explain its role: ensuring $\|\mathsf{PE}(p)\| = 1$ and making the encoding kernel converge to $k(\Delta)$ rather than diverge to $O(d)$ as $d\to\infty$.

\textbf{2b. $L$-dependence removed from boxed formulas}:

\begin{longtable}[]{@{}
  >{\raggedright\arraybackslash}p{(\linewidth - 6\tabcolsep) * \real{0.2308}}
  >{\raggedright\arraybackslash}p{(\linewidth - 6\tabcolsep) * \real{0.2308}}
  >{\raggedright\arraybackslash}p{(\linewidth - 6\tabcolsep) * \real{0.2308}}
  >{\raggedright\arraybackslash}p{(\linewidth - 6\tabcolsep) * \real{0.3077}}@{}}
\toprule\noalign{}
\begin{minipage}[b]{\linewidth}\raggedright
File
\end{minipage} & \begin{minipage}[b]{\linewidth}\raggedright
Location
\end{minipage} & \begin{minipage}[b]{\linewidth}\raggedright
Formula
\end{minipage} & \begin{minipage}[b]{\linewidth}\raggedright
$L$ dependence
\end{minipage} \\
\midrule\noalign{}
\endhead
\bottomrule\noalign{}
\endlastfoot
main.tex & L317--320 & $\lambda_j = 2\pi\xi \cdot \cot(\pi(2j+1)/2d)$ & None \checkmark \\
ppe\_rigorous\_derivation.md & L133 & $\lambda_j = 2\pi\xi \cdot \cot(\pi(2j+1)/2d)$ & None \checkmark \\
\end{longtable}

$L$ only enters indirectly through Corollary 1.2.1's Nyquist condition $d_{\min} \approx L/(2\xi)$, not in the encoding's boxed formula itself. \textbf{PASS.}

\begin{center}\rule{0.5\linewidth}{0.5pt}\end{center}

\subsubsection{3. Theorem Numbering Unified to 4.1?}\label{thm-numbering-unified}

``Theorem 4.1'' refers to the theorem on ``precise conditions under which learned PE breaks Theorem 3.''

\begin{longtable}[]{@{}llll@{}}
\toprule\noalign{}
File & Location & Numbering & Label \\
\midrule\noalign{}
\endhead
\bottomrule\noalign{}
\endlastfoot
main.tex & L1095 & Theorem 4.1 & \texttt{{\textbackslash}label\{thm:4.1\}} \\
ppe\_rigorous\_derivation.md & L775 & Theorem 4.1 & --- \\
\end{longtable}

Both files uniformly use \textbf{Theorem 4.1}. Proposition 4.1 (fixed PE preserves indistinguishability) is also consistently numbered. \textbf{PASS.}

\begin{center}\rule{0.5\linewidth}{0.5pt}\end{center}

\subsubsection{4. $d_{\min}$ Corrected to $L/(2\xi)$?}\label{dmin-corrected}

\begin{longtable}[]{@{}
  >{\raggedright\arraybackslash}p{(\linewidth - 6\tabcolsep) * \real{0.2500}}
  >{\raggedright\arraybackslash}p{(\linewidth - 6\tabcolsep) * \real{0.2500}}
  >{\raggedright\arraybackslash}p{(\linewidth - 6\tabcolsep) * \real{0.2500}}
  >{\raggedright\arraybackslash}p{(\linewidth - 6\tabcolsep) * \real{0.2500}}@{}}
\toprule\noalign{}
\begin{minipage}[b]{\linewidth}\raggedright
File
\end{minipage} & \begin{minipage}[b]{\linewidth}\raggedright
Location
\end{minipage} & \begin{minipage}[b]{\linewidth}\raggedright
Formula
\end{minipage} & \begin{minipage}[b]{\linewidth}\raggedright
Derivation
\end{minipage} \\
\midrule\noalign{}
\endhead
\bottomrule\noalign{}
\endlastfoot
main.tex & L368 & $d_{\min} \approx \frac{L}{2\xi}$ & $\lambda_{\max} \approx 4d\xi$, Nyquist: $\lambda_{\max} \geq 2L \to 4d\xi \geq 2L \to d \geq L/(2\xi)$ \\
ppe\_rigorous\_derivation.md & L173 & $d_{\min} \approx \frac{L}{2\xi}$ & Same \\
\end{longtable}

Derivation chain complete, Nyquist condition correctly cited (covering the entire sequence interval requires $\lambda_{\max} \geq 2L$). \textbf{PASS.}

\begin{center}\rule{0.5\linewidth}{0.5pt}\end{center}

\subsubsection{5. Theorem 2.4.1 Downgraded to Heuristic Estimate?}\label{thm-2.4.1-downgraded}

\textbf{Key change}: From ``Theorem (lower bound)'' $\to$ ``Proposition (heuristic estimate --- not a rigorous lower bound).''

\begin{longtable}[]{@{}
  >{\raggedright\arraybackslash}p{(\linewidth - 4\tabcolsep) * \real{0.1395}}
  >{\raggedright\arraybackslash}p{(\linewidth - 4\tabcolsep) * \real{0.2326}}
  >{\raggedright\arraybackslash}p{(\linewidth - 4\tabcolsep) * \real{0.6279}}@{}}
\toprule\noalign{}
\begin{minipage}[b]{\linewidth}\raggedright
Feature
\end{minipage} & \begin{minipage}[b]{\linewidth}\raggedright
main.tex
\end{minipage} & \begin{minipage}[b]{\linewidth}\raggedright
ppe\_rigorous\_derivation.md
\end{minipage} \\
\midrule\noalign{}
\endhead
\bottomrule\noalign{}
\endlastfoot
Environment type & \texttt{{\textbackslash}begin\{proposition\}} & \texttt{**Proposition 2.4.1**} \\
Honesty label & \texttt{{\textbackslash}heuristic\{\}} + warning box & \texttt{[Heuristic]} + warning box \\
Explicit statement & ``This expression is \textbf{not} a rigorous mathematical lower bound'' & ``This expression is \emph{not} a rigorous mathematical lower bound'' \\
Logical error noted & ``One cannot logically deduce a lower bound for the difference from two Fano lower bounds'' & ``From two Fano lower bounds\ldots cannot logically deduce'' \\
Equality misuse explanation & ``If the Fano inequality attains equality in both worlds (which generally does not hold)'' & ``If one assumes Fano inequality attains equality in both worlds (requires specific distribution conditions)'' \\
\end{longtable}

Both files well aligned. \textbf{PASS.}

\begin{center}\rule{0.5\linewidth}{0.5pt}\end{center}

\subsubsection{6. Theorem 2.2.1 Step 4 Restricted to Bayes Optimal?}\label{thm-2.2.1-step4-restricted}

\begin{longtable}[]{@{}
  >{\raggedright\arraybackslash}p{(\linewidth - 4\tabcolsep) * \real{0.1395}}
  >{\raggedright\arraybackslash}p{(\linewidth - 4\tabcolsep) * \real{0.2326}}
  >{\raggedright\arraybackslash}p{(\linewidth - 4\tabcolsep) * \real{0.6279}}@{}}
\toprule\noalign{}
\begin{minipage}[b]{\linewidth}\raggedright
Aspect
\end{minipage} & \begin{minipage}[b]{\linewidth}\raggedright
main.tex
\end{minipage} & \begin{minipage}[b]{\linewidth}\raggedright
ppe\_rigorous\_derivation.md
\end{minipage} \\
\midrule\noalign{}
\endhead
\bottomrule\noalign{}
\endlastfoot
Theorem title & ``Sufficient condition --- information-theoretic form (Bayes optimal restricted)'' & ``Sufficient condition --- information-theoretic form, Bayes optimal restricted'' \\
Conclusion statement & ``\textbf{Bayes optimal classifier} in augmented features satisfies $\delta>0$'' & ``\textbf{Bayes optimal classifier} in augmented features satisfies $\delta>0$'' \\
Step 4 marking & \texttt{{\textbackslash}heuristic\{\}} + ``conjecture'' & \texttt{[Heuristic]} + ``conjecture'' \\
Limitation statement & ``Steps 1--3 are rigorous information-theoretic derivations (conclusion holds for Bayes optimal classifier)'' & ``Steps 1--3 prove the Bayes optimal classifier can benefit from position information'' \\
Degradation statement & ``If the generalization of Step 4 fails, the theorem degrades to a classical conclusion --- trivial'' & ``If the generalization of Step 4 fails, Theorem 2.2.1 degrades to a classical conclusion --- trivial'' \\
\end{longtable}

\textbf{PASS.}

\begin{center}\rule{0.5\linewidth}{0.5pt}\end{center}

\subsubsection{7. Theorem 3' Breakage Direction Changed to Open Problem?}\label{thm-3prime-open-problem}

\begin{longtable}[]{@{}
  >{\raggedright\arraybackslash}p{(\linewidth - 4\tabcolsep) * \real{0.1395}}
  >{\raggedright\arraybackslash}p{(\linewidth - 4\tabcolsep) * \real{0.2326}}
  >{\raggedright\arraybackslash}p{(\linewidth - 4\tabcolsep) * \real{0.6279}}@{}}
\toprule\noalign{}
\begin{minipage}[b]{\linewidth}\raggedright
Location
\end{minipage} & \begin{minipage}[b]{\linewidth}\raggedright
main.tex
\end{minipage} & \begin{minipage}[b]{\linewidth}\raggedright
ppe\_rigorous\_derivation.md
\end{minipage} \\
\midrule\noalign{}
\endhead
\bottomrule\noalign{}
\endlastfoot
Within Theorem 3' statement & ``If premise (A) fails: whether indistinguishability is broken is currently an \textbf{open problem}'' & ``If premise (A) fails: whether indistinguishability is broken is currently an \textbf{open problem}'' \\
Honesty label & \texttt{{\textbackslash}openproblem\{\}} & \texttt{[Open Problem]} \\
Minimal example reassessment & ``This example actually proves Theorem 3's robustness, not its fragility'' & ``This minimal example actually proves Theorem 3's robustness, not its fragility'' \\
Honesty table & Theorem 4.1 entry: \texttt{{\textbackslash}openproblem\ breakage direction lacks constructive proof} & --- \\
Theorem 3' entry & \texttt{{\textbackslash}rigorous\ when premise holds;\ {\textbackslash}openproblem\ when premise fails} & --- \\
\end{longtable}

The originally claimed ``constructive example proves breakage direction'' has been reassessed as ``the constructive example instead demonstrates the theorem's robustness'' --- as expected. \textbf{PASS.}

\begin{center}\rule{0.5\linewidth}{0.5pt}\end{center}

\subsubsection{8. Any Remaining Symbol Inconsistencies Between the Two Files?}\label{remaining-symbol-inconsistencies}

\paragraph{8a. Consistent items \checkmark}\label{consistent-items}

\begin{longtable}[]{@{}
  >{\raggedright\arraybackslash}p{(\linewidth - 6\tabcolsep) * \real{0.2037}}
  >{\raggedright\arraybackslash}p{(\linewidth - 6\tabcolsep) * \real{0.1852}}
  >{\raggedright\arraybackslash}p{(\linewidth - 6\tabcolsep) * \real{0.5000}}
  >{\raggedright\arraybackslash}p{(\linewidth - 6\tabcolsep) * \real{0.1111}}@{}}
\toprule\noalign{}
\begin{minipage}[b]{\linewidth}\raggedright
Symbol/Concept
\end{minipage} & \begin{minipage}[b]{\linewidth}\raggedright
main.tex
\end{minipage} & \begin{minipage}[b]{\linewidth}\raggedright
ppe\_rigorous\_derivation.md
\end{minipage} & \begin{minipage}[b]{\linewidth}\raggedright
Status
\end{minipage} \\
\midrule\noalign{}
\endhead
\bottomrule\noalign{}
\endlastfoot
Detection margin & $\Delta_s = p_{\text{noisy}} - p_{\text{clean}}$ & $\Delta_s = p_{\text{noisy}} - p_{\text{clean}}$ & \checkmark \\
PPE augmented margin & $\Delta_s^{\text{Situs}} = \Delta_s + \delta_s^{\text{PE}}$ & $\Delta_s^{\text{PPE}} = \Delta_s + \delta_s^{\text{PE}}$ & \checkmark \\
Chernoff-Hoeffding & $\exp(-2M(\Delta_s + \delta_s^{\text{PE}})^2)$ & $\exp(-2M(\Delta_s + \delta_s^{\text{PE}})^2)$ & \checkmark \\
Encoding imperfection & $\varepsilon_{\text{PE}} = I(Y;P|X) - I(Y;\text{PE}(P)|X)$ & $\varepsilon_{\text{PE}} = I(Y;P|X) - I(Y;\text{PE}(P)|X)$ & \checkmark \\
Theorem 2' upper bound & $F1_{\text{base}} + C_F\cdot\sqrt{(\delta + 2\varepsilon/C_F^2)/2}$ & $F1_{\text{base}} + C_F\cdot\sqrt{(\delta + 2\varepsilon/C_F^2)/2}$ & \checkmark \\
$d_{\min}$ & $L/(2\xi)$ & $L/(2\xi)$ & \checkmark \\
Scalar PE definition & $\sqrt{2/d}\cdot\sin/\cos(2\pi p/\lambda_j)$ & $\sqrt{2/d}\cdot\sin/\cos(2\pi p/\lambda_j)$ & \checkmark \\
3D rotation PE & $\mathbf{R}(\mathbf{p})\cdot\mathbf{e}_0$ & $\mathbf{R}(\mathbf{p})\cdot\mathbf{e}_0$ & \checkmark \\
Rotation Lipschitz & $\max(\alpha, \beta, \gamma)$ & $\max(\alpha, \beta, \gamma)$ & \checkmark \\
Sine Lipschitz & $(2\sqrt{2}\cdot\pi/\sqrt{d})\cdot\sqrt(\sum 1/\lambda_j^2)$ & $(2\sqrt{2}\cdot\pi/\sqrt{d})\cdot\sqrt(\sum 1/\lambda_j^2)$ & \checkmark \\
Theorem 4.1 condition & $\mathcal{L}_{\text{total}} = \mathcal{L}_{\text{sup}} + \lambda\cdot\mathcal{L}_{\text{aux}}$ & $\mathcal{L}_{\text{total}} = \mathcal{L}_{\text{sup}} + \lambda\cdot\mathcal{L}_{\text{aux}}$ & \checkmark \\
\end{longtable}

\paragraph{8b. Expected differences (non-issues)}\label{expected-differences}

\begin{longtable}[]{@{}
  >{\raggedright\arraybackslash}p{(\linewidth - 4\tabcolsep) * \real{0.3333}}
  >{\raggedright\arraybackslash}p{(\linewidth - 4\tabcolsep) * \real{0.3333}}
  >{\raggedright\arraybackslash}p{(\linewidth - 4\tabcolsep) * \real{0.3333}}@{}}
\toprule\noalign{}
\begin{minipage}[b]{\linewidth}\raggedright
Difference
\end{minipage} & \begin{minipage}[b]{\linewidth}\raggedright
Reason
\end{minipage} & \begin{minipage}[b]{\linewidth}\raggedright
Judgment
\end{minipage} \\
\midrule\noalign{}
\endhead
\bottomrule\noalign{}
\endlastfoot
\texttt{{\textbackslash}Situs\{\}} vs \texttt{PPE} & paper uses framework name Situs, derivation uses abbreviation PPE & Expected \checkmark \\
$\hat{\ }\{{\textbackslash}\text{Situs}\}$ vs $\hat{\ }\{{\textbackslash}\text{PPE}\}$ & Different superscript conventions & Expected \checkmark \\
Section numbering offset & Paper has \S1 introduction, derivation starts directly from encoding definitions without introduction & Expected \checkmark \\
Theorem environment format & LaTeX \texttt{{\textbackslash}begin\{theorem\}} vs Markdown \texttt{**Theorem**} & Format difference \checkmark \\
\texttt{{\textbackslash}PE} vs \texttt{PE} & LaTeX operator vs plain text & Format difference \checkmark \\
\end{longtable}

\paragraph{8c. \warning\ Residual inconsistencies (require attention)}\label{residual-inconsistencies}

\textbf{R1. ``Necessary and sufficient condition'' overclaim in Abstract}

\begin{longtable}[]{@{}
  >{\raggedright\arraybackslash}p{(\linewidth - 2\tabcolsep) * \real{0.5000}}
  >{\raggedright\arraybackslash}p{(\linewidth - 2\tabcolsep) * \real{0.5000}}@{}}
\toprule\noalign{}
\begin{minipage}[b]{\linewidth}\raggedright
Location
\end{minipage} & \begin{minipage}[b]{\linewidth}\raggedright
Content
\end{minipage} \\
\midrule\noalign{}
\endhead
\bottomrule\noalign{}
\endlastfoot
main.tex L97 (Abstract) & ``proves the \textbf{necessary and sufficient condition} for $\delta_s^{\text{PE}} > 0$ is $I(Y; P \mid S) > 0$ (Theorem 2.2.1)'' \\
main.tex L649 (Theorem 2.2.1 title) & ``\textbf{Sufficient condition} for $\delta_s^{\text{PE}} > 0$ --- information-theoretic form (Bayes optimal restricted)'' \\
main.tex L694 (Proposition 2.2.1 title) & ``\textbf{Necessary condition} for $\delta_s^{\text{PE}} > 0$'' \\
ppe\_rigorous\_derivation.md L1020 (Summary) & ``proved \textbf{sufficient condition} ($I(Y; P \mid S) > 0$) for $\delta_s^{\text{PE}} > 0$'' \\
\end{longtable}

\textbf{Problem}: The abstract claims Theorem 2.2.1 gives a ``necessary and sufficient condition,'' but the actual Theorem 2.2.1 only proves a ``sufficient condition'' (and is restricted to the Bayes optimal classifier), while the necessary condition is separately given by Proposition 2.2.1 (and allows non-information-theoretic mechanisms). Moreover, the sufficient condition involves the heuristic generalization of Step 4 (from Bayes optimal to actual experts), not a purely rigorous proof.

\textbf{Suggestion}: The abstract's ``necessary and sufficient condition'' should be changed to ``sufficient condition (Theorem 2.2.1) and necessary condition (Proposition 2.2.1),'' or at least the attribution should be corrected from ``Theorem 2.2.1'' alone to ``Theorem 2.2.1 + Proposition 2.2.1.''

\textbf{R2. ppe\_rigorous\_derivation.md Theorem 2.3.1's dual statement}

\begin{longtable}[]{@{}
  >{\raggedright\arraybackslash}p{(\linewidth - 2\tabcolsep) * \real{0.5000}}
  >{\raggedright\arraybackslash}p{(\linewidth - 2\tabcolsep) * \real{0.5000}}@{}}
\toprule\noalign{}
\begin{minipage}[b]{\linewidth}\raggedright
Location
\end{minipage} & \begin{minipage}[b]{\linewidth}\raggedright
Content
\end{minipage} \\
\midrule\noalign{}
\endhead
\bottomrule\noalign{}
\endlastfoot
L458 & Rough form: $\min(1, \sqrt{\frac12\cdot I(P;\text{PE}(P)|S,Y)}) + \text{Bernstein correction}$ \\
L462 & Precise form: $\min(1-p_{\text{clean}}, \sqrt{\frac12\cdot D_{\text{KL}}(P_{\text{PE}|S,Y} \| P_{\text{PE}|S})}) + \delta_s^{\text{variance}}$ \\
main.tex L723 & Only precise form \checkmark \\
\end{longtable}

\textbf{Problem}: ppe\_rigorous\_derivation.md first gives a rougher form, then a precise form. The two versions differ in their second terms ($I(P;\text{PE}|S,Y)$ vs $D_{\text{KL}}$), and the first version uses a conceptually different ``Bernstein correction'' while the precise version uses $\delta_s^{\text{variance}}$. This may confuse readers --- the two forms are not equivalent.

\textbf{Suggestion}: ppe\_rigorous\_derivation.md should add a ``roughly'' label before the rough form and explain that the precise form follows below. Or directly delete the rough form, keeping only the precise version consistent with main.tex.

\begin{center}\rule{0.5\linewidth}{0.5pt}\end{center}

\subsection{Overall Assessment}\label{overall-assessment}

\begin{longtable}[]{@{}ll@{}}
\toprule\noalign{}
Check Item & Status \\
\midrule\noalign{}
\endhead
\bottomrule\noalign{}
\endlastfoot
1. Lipschitz constant corrected & \checkmark\ PASS \\
2. Normalization factor + $L$-dependence removed & \checkmark\ PASS \\
3. Theorem numbering unified to 4.1 & \checkmark\ PASS \\
4. $d_{\min} = L/(2\xi)$ & \checkmark\ PASS \\
5. Theorem 2.4.1 downgraded to heuristic & \checkmark\ PASS \\
6. Theorem 2.2.1 Step 4 Bayes optimal restricted & \checkmark\ PASS \\
7. Theorem 3' breakage direction $\to$ open problem & \checkmark\ PASS \\
8. Symbol consistency between two files & \warning\ 2 residual issues (see R1, R2) \\
\end{longtable}

\textbf{Summary}: 7/8 items fully pass. Core corrections (Lipschitz constant, normalization factor, Fano downgrade, $d_{\min}$ correction, Bayes optimal restriction, breakage direction open-problem-ization) have been correctly implemented and are consistent across both files. The 2 residual issues are minor level --- Abstract wording is overly absolute, and the derivation document has dual theorem statements --- neither affects mathematical correctness and can be corrected in the next editing round.

\end{document}
